En este último capítulo se detallan las lecciones aprendidas tras el desarrollo del presente proyecto y se identifican las posibles oportunidades de mejora sobre el \textit{software} desarrollado.

\section{Objetivos alcanzados}

El objetivo principal de este Trabajo Fin de Grado era desarrollar una aplicación móvil multiplataforma que permitiese consultar, visualizar y analizar información futbolística mediante técnicas de inteligencia de negocio, minería de datos y aprendizaje automático. Para ello, se planteó una solución que integrase datos históricos de LaLiga, una arquitectura de almacenamiento orientada al análisis, una API de consulta, una interfaz móvil y una capa de análisis avanzado.

En primer lugar, se ha conseguido construir una base de datos estructurada a partir de información histórica de la Primera División Española. Para ello, se ha desarrollado un proceso de extracción, transformación y carga de datos que permite pasar de información procedente de una API externa a un modelo físico organizado y preparado para su explotación. Este proceso ha permitido disponer de datos relacionados con temporadas, equipos, jugadores, partidos, eventos y estadísticas de rendimiento.

También se ha alcanzado el objetivo de diseñar una arquitectura de inteligencia de negocio basada en un modelo de datos orientado al análisis. La base de datos se ha organizado mediante un esquema en constelación, formado por tablas de dimensiones y tablas de hechos, lo que facilita la consulta de información desde diferentes perspectivas. Gracias a esta estructura, la aplicación puede mostrar datos a nivel de temporada, equipo, jugador y partido de forma ordenada y coherente.

Otro de los objetivos alcanzados ha sido la implementación de una API REST que actúa como capa intermedia entre la aplicación móvil y la base de datos. Esta API permite separar la lógica de acceso a datos de la capa visual, facilitando la consulta de información desde las distintas pantallas de la aplicación. A través de ella se obtienen listados, detalles, rankings, gráficos, favoritos, predicciones y respuestas asociadas al asistente conversacional.

En cuanto a la aplicación móvil, se ha desarrollado una interfaz que permite navegar entre las secciones principales del sistema: inicio, temporadas, equipos, jugadores, partidos, favoritos, perfil, simulación, análisis mediante \textit{data mining} y asistente inteligente. Las pantallas de detalle y los dashboards permiten representar la información de manera visual e interactiva, evitando que el usuario tenga que interpretar directamente tablas extensas o datos sin procesar.

Además de la consulta estadística tradicional, se han incorporado resultados generados mediante técnicas de minería de datos y aprendizaje automático. Entre ellos se incluyen predicciones de partidos, simulaciones de clasificación, estados de forma de equipos y jugadores, jugadores similares, probables goleadores y recomendaciones de fichajes. Estos resultados se han centralizado en una pantalla específica de \textit{data mining}, dando mayor protagonismo a la parte analítica del proyecto.

Por último, se ha integrado un asistente conversacional basado en inteligencia artificial, denominado Agente de Oro. Esta funcionalidad permite al usuario realizar preguntas en lenguaje natural sobre la información almacenada en el sistema, facilitando el acceso a datos concretos sin necesidad de recorrer manualmente todas las pantallas de la aplicación.

En conjunto, puede afirmarse que el proyecto ha cumplido con su objetivo general: desarrollar una aplicación móvil capaz de acercar el análisis avanzado de datos futbolísticos a usuarios no necesariamente especializados, combinando visualización, predicción, simulación y consulta inteligente dentro de una misma plataforma.

\section{Lecciones aprendidas}

El desarrollo del proyecto ha permitido extraer varias conclusiones relevantes tanto desde el punto de vista técnico como desde el punto de vista funcional.

Una de las principales lecciones aprendidas es la importancia de contar con una estructura de datos sólida antes de desarrollar las funcionalidades visibles de la aplicación. En un proyecto basado en análisis deportivo, la calidad de los resultados depende directamente de la calidad, coherencia y organización de los datos disponibles. Por este motivo, el diseño del modelo de datos y el proceso de transformación han sido elementos fundamentales para que posteriormente pudieran construirse dashboards, rankings y modelos predictivos.

También se ha comprobado que la visualización de datos es tan importante como el propio cálculo de las métricas. Muchos indicadores deportivos pueden resultar complejos si se presentan únicamente en forma de tablas o valores numéricos. Sin embargo, al representarlos mediante tarjetas, gráficos, rankings y dashboards, la información resulta mucho más comprensible para el usuario. Esta idea ha sido especialmente relevante en las pantallas de temporada, equipo, jugador, partido y \textit{data mining}.

Otra lección importante es que los modelos de minería de datos aportan valor cuando sus resultados se integran de forma clara dentro de la aplicación. No basta con calcular predicciones, simulaciones o agrupaciones de jugadores; es necesario presentarlas de una manera que el usuario pueda entender e interpretar. Por ello, la creación de una pantalla específica de \textit{data mining} ha permitido dar mayor visibilidad a estos resultados y diferenciarlos de la consulta estadística tradicional.

El desarrollo también ha puesto de manifiesto la utilidad de separar correctamente las distintas capas del sistema. La división entre aplicación móvil, backend, base de datos y procesos analíticos facilita el mantenimiento del proyecto y permite modificar una parte del sistema sin afectar directamente al resto. Esta separación ha sido especialmente útil al añadir nuevas funcionalidades, como favoritos, simulaciones, predicciones o el asistente conversacional.

En relación con el asistente inteligente, se ha observado que el uso de lenguaje natural puede mejorar la experiencia de usuario, especialmente cuando se desea consultar información concreta. No obstante, también se ha comprobado que esta funcionalidad requiere restricciones claras para evitar consultas incorrectas, ambiguas o fuera del dominio del sistema. Por ello, ha sido necesario limitar su uso a consultas relacionadas con la información futbolística almacenada y controlar que las operaciones generadas sean seguras.

Por último, el proyecto ha servido para comprender que una aplicación de análisis deportivo no debe centrarse únicamente en acumular funcionalidades, sino en organizar la información de forma útil. La prioridad debe ser que el usuario pueda encontrar rápidamente lo que busca, interpretar los resultados y entender el valor de los datos mostrados.

\section{Trabajo futuro}

Aunque el proyecto ha alcanzado los objetivos planteados, existen diferentes líneas de trabajo futuro que permitirían ampliar sus capacidades y acercarlo a un entorno de uso real.

Una primera mejora sería ampliar la cobertura de datos. Actualmente, el sistema se centra en la Primera División Española, por lo que una evolución natural consistiría en incorporar otras competiciones, como Segunda División, competiciones europeas o ligas extranjeras. Esto permitiría comparar equipos y jugadores entre distintos contextos y enriquecer los análisis disponibles.

Otra línea de mejora sería incorporar datos en tiempo real. En el estado actual del proyecto, la aplicación trabaja principalmente con datos históricos y datos previamente cargados. En un entorno de producción, sería interesante integrar servicios que proporcionen información en directo sobre partidos, eventos, alineaciones, lesiones o estadísticas actualizadas durante el encuentro. Esto permitiría ofrecer predicciones y análisis más dinámicos.

También sería posible mejorar los modelos de predicción mediante la incorporación de nuevas variables. Por ejemplo, podrían añadirse datos sobre lesiones, sanciones, alineaciones probables, calendario acumulado, descanso entre partidos, mercado de fichajes o métricas avanzadas como xG y xA cuando estuvieran disponibles. Estos factores podrían aumentar la precisión de los modelos y aportar una visión más completa del rendimiento deportivo.

En relación con la simulación de temporada, una posible mejora consistiría en permitir al usuario crear y guardar múltiples escenarios personalizados. De esta forma, podría comparar distintas hipótesis sobre resultados futuros y analizar cómo afectarían a la clasificación final, a la lucha por el título, a la clasificación europea o al descenso.

La pantalla de \textit{data mining} también podría evolucionar hacia una sección más personalizada. Actualmente centraliza los resultados analíticos del sistema, pero en futuras versiones podría priorizar automáticamente la información relacionada con los equipos y jugadores favoritos del usuario. Esto permitiría ofrecer una experiencia más adaptada a los intereses de cada persona.

Otra línea de trabajo relevante sería mejorar el asistente conversacional. En futuras versiones, el sistema podría ser capaz de interpretar preguntas más complejas, explicar gráficos, comparar jugadores o equipos y generar resúmenes automáticos de partidos, temporadas o trayectorias. También sería interesante mejorar el control de las consultas generadas para aumentar la seguridad y fiabilidad de las respuestas.

Desde el punto de vista de despliegue, el proyecto podría evolucionar hacia una versión en producción. Para ello sería necesario alojar el backend y la base de datos en un servidor externo, configurar un dominio, establecer copias de seguridad, monitorizar el sistema y valorar la publicación de la aplicación en tiendas móviles. Además, habría que considerar el coste recurrente de servicios externos, como APIs deportivas o modelos de inteligencia artificial.

Finalmente, sería recomendable realizar pruebas con usuarios reales. Aunque durante el desarrollo se han realizado pruebas funcionales y de integración, una validación con aficionados, analistas deportivos o perfiles técnicos permitiría obtener información más precisa sobre la utilidad de la aplicación, la claridad de los dashboards, la comprensión de las predicciones y la facilidad de uso del asistente conversacional.

En definitiva, el trabajo realizado constituye una base funcional sobre la que pueden desarrollarse nuevas mejoras. La aplicación demuestra que es posible integrar datos deportivos, visualización, modelos predictivos y consultas inteligentes en una herramienta móvil accesible, quedando abiertas varias líneas de evolución hacia un sistema más completo, personalizado y cercano a un uso profesional.